% Abstract: one paragraph, under 250 words, no citations. Processed from the
% draft instruction: workflow overview, headline metrics, the four-entrant
% competition with the second-place explanation, and the external-standards
% strength. File names kept short and only where they carry weight.

In this study, an autonomous Claude Code agent carried an instruction specification through a generated codebase and full
execution record regarding a 2030 Unitree
H2-Surgical 1.0 humanoid performing a 60-second pancreaticoduodenectomy. The surgical humanoid code assurance process, not code
generation, is the substantial and decision-bearing work; holding verification,
validation, and uncertainty quantification to a higher standard than the code
itself.\ This process is what will make physical AI oncology trials faster, less
expensive, and more beneficial to patients.\ The run passed 172 of 172 automated
tests, 64 of them in the ten-gate assurance suite. Each gate required a
verification fraction of exactly 1.0, validation agreement up to 1.00, relative
error as tight as 0.01, and a coefficient-of-variation bound as tight as 0.05;
five decision cases resolved to ten ACCEPT, three BLOCK, and one ESCALATE, and
all 32 of 32 sweep iterations cleared every gate near a composite mean of 93.6.
In a four-entrant, simulation-against-simulation tournament the mobile humanoid
placed second (93.334) to the stationary eight-arm PancreSpeed cart (93.782). Parallel arms
shorten the throughput-weighted score, so the single humanoid trails by under
half a point, whereas the prior PDAC paper featured the multi-arm baseline alone.
The 1790-line \texttt{comparison.json} and the 1000-row
\texttt{sample\_h2\_sensor.csv} reproduced byte-for-byte from seed 20260525.
Binding every gate to published consensus standards already used in real life
makes the assurance argument traceable and defensible to a regulator.
